% ===========================================================================
% Section II-A -- Continuous-to-Discrete Fault Signal Model
% Owner: Member 1
% ===========================================================================
\subsection{Continuous-to-Discrete Fault Signal Model}
A localized bearing defect behaves like a sequence of short impacts. Write the
impact times as $t_m = m/f_{\text{fault}} + \varepsilon_m$, where
$f_{\text{fault}}$ is the nominal repetition rate and $\varepsilon_m$ is a
timing error from local slip or speed fluctuation. Impact $m$ has amplitude
$A_m$ and excites the structure with impulse response $h(t)$. Background
machine vibration and sensor noise enter as $w(t)$, so
\begin{equation}
  x(t) = \sum_{m} A_m\, h(t-t_m) + w(t).
  \label{eq:ctmodel}
\end{equation}
For an underdamped structural mode the causal impulse response is a decaying
sinusoid,
\begin{equation}
  h(t) = e^{-\sigma t}\cos(2\pi f_r t)\,u(t), \qquad t\geq 0,
  \label{eq:ctir}
\end{equation}
where $\sigma$ is the decay rate, $f_r$ the resonance frequency and $u(t)$ the
unit step. The response starts at the impact and decays afterwards, which is
what makes it causal. The same repeated-impact, exponentially decaying picture
is the classical localized-defect model~\cite{mcfadden1984}.

Sampling with period $T_s=1/f_s$ gives
\begin{equation}
  x[n] = \sum_m A_m\, h[n-n_m] + w[n], \qquad
  n_m \approx \operatorname{round}(f_s t_m).
  \label{eq:model}
\end{equation}
The same signal can be written through the impact sequence
$p[n]=\sum_m A_m\delta[n-n_m]$, as $x[n]=p[n]*h[n]+w[n]$, where $*$ is linear
convolution. The nominal spacing between impacts is
$N_p\approx f_s/f_{\text{fault}}$ samples. $N_p$ need not be an integer, and
small timing errors push the observed spectral peaks away from the ideal
kinematic frequencies.

The sampling frequency has to be more than twice the highest vibration
frequency that must be represented, $f_s>2f_{\max}$. The acquisition chain also
has to suppress out-of-band components before sampling, or aliasing folds them
into the band of interest. The records used here run at $12$ and $48$~kHz
(CWRU) and at $48.828$ and $97.656$~kHz (MFPT), which puts the Nyquist
frequencies at $6$, $24$, $24.414$ and $48.828$~kHz.

In the $z$-plane, $h[n]$ gives a complex-conjugate pole pair
\begin{equation}
  z_{1,2} = r e^{\pm j\omega_r}, \qquad r = e^{-\sigma T_s} < 1,
  \label{eq:poles}
\end{equation}
and stability follows from $|z_{1,2}|<1$. The model also makes the central
point of the paper explicit. The fault signature $A_m$ is a
\emph{low-frequency modulation} of a \emph{high-frequency carrier}
$\omega_r$. Demodulation separates the two; direct spectral inspection does
not.
